\subsection{Symmetry and Self-adjointness}\label{sec:symmetry-and-self-adjointness}\label{sec:adjoint-operators}

\subsubsection{The Adjoint Operator}\label{subsec:banach-adjoint-operator}

We now define the \emph{adjoint} of an operator.

\begin{Definition}
\label{def:banach-adjoint-operator}
Let $(\mathcal{D}(A),A)$ be an operator from a Banach space $X$ to a Banach space $Y$ such that $\mathcal{D}(A)$ is dense in $X$. We define $\mathcal{D}(A^\times)$ to be the set of all $v\in Y^*$ for which there exists $w\in X^*$ such that
\begin{equation}\label{eq:banach-adjoint-dual-pairing-definition}
v(Au)=w(u)
\end{equation}
for all $u\in\mathcal{D}(A)$. Note that since $\mathcal{D}(A)$ is dense, $w$ is uniquely determined by $v\in\mathcal{D}(A^\times)$ and \eqref{eq:banach-adjoint-dual-pairing-definition}. Thus, we can define an operator $(\mathcal{D}(A^\times),A^\times)$ from $Y^*$ to $X^*$ by
\begin{equation}\label{eq:banach-adjoint-action-definition}
A^\times(v):=w
\end{equation}
for every $v\in\mathcal{D}(A^\times)$. We call $(\mathcal{D}(A^\times),A^\times)$ the \textbf{adjoint} of $(\mathcal{D}(A),A)$.
\end{Definition}

It is clear that $\mathcal{D}(A^\times)$ is nonempty since $\{0\}\in\mathcal{D}(A^\times)$. Also, it follows directly from the definition that $A^\times$ is linear. Furthermore, for bounded operators we can show the following.

\begin{Theorem}
\label{thm:bounded-adjoint-domain-and-norm}
For any bounded operator $A\in\mathcal{L}(X,Y)$ we have $\mathcal{D}(A^\times)=Y^*$ and $A^\times:Y^*\to X^*$ is a bounded operator with $\|A^\times\|=\|A\|$.
\end{Theorem}

The proof depends on the following lemma, which is a direct consequence of the Hahn-Banach theorem (Theorem~\ref{thm:hahn-banach-extension}).

\begin{Lemma}
\label{lem:unit-norm-norming-functional}
Let $X$ be a Banach space and let $\overline x$ be any nonzero element of $X$. Then there exists a linear functional $l\in X^*$ such that
\begin{equation}\label{eq:norming-functional-unit-norm-equality}
\|l\|=1\quad\text{and}\quad l(\overline x)=\|\overline x\|.
\end{equation}

\end{Lemma}

\begingroup
\postdisplaypenalty=10000
\renewcommand{\endProof}{\par\nobreak\hfill$\blacksquare$\par\vspace{0.1cm}}
\begin{Proof}
Let $M:=\{\alpha\overline x\mid\alpha\in\mathbb{R}\}$ be the subspace spanned by $\overline x$. We define a linear functional $\widetilde l$ on $M$ by
\begin{equation}\label{eq:norming-functional-one-dimensional-definition}
\widetilde l(\alpha\overline x)=\alpha\|\overline x\|.
\end{equation}
It is easy to see that $\widetilde l$ has norm $1$. The Hahn-Banach theorem assures us that $\widetilde l$ has an extension $l$ to all of $X$ with norm less than or equal $1$. Since $l(\overline x)=\widetilde l(\overline x)=\|\overline x\|$ we see that in fact the norm is equal to $1$, and the lemma is proved.
\end{Proof}
\endgroup

We now prove Theorem~\ref{thm:bounded-adjoint-domain-and-norm}.

\begingroup
\postdisplaypenalty=10000
\renewcommand{\endProof}{\par\nobreak\hfill$\blacksquare$\par\vspace{0.1cm}}
\begin{Proof}
For any bounded linear functional $v\in Y^*$ we see that
\begin{equation}\label{eq:bounded-adjoint-induced-functional}
u\mapsto v(A(u))=:w(u)
\end{equation}
is a linear map from $X$ to $\mathbb{R}$. We further see that this map is bounded since
\begin{equation}\label{eq:bounded-adjoint-functional-estimate}
|w(u)|=|v(A(u))|\leq\|v\|\|A(u)\|\leq\|v\|\|A\|\|u\|.
\end{equation}
Thus, $v\in\mathcal{D}(A^\times)$ and $w=A^\times(v)$. We can also get from \eqref{eq:bounded-adjoint-functional-estimate} that
\begin{equation}\label{eq:bounded-adjoint-image-norm-estimate}
\|A^\times(v)\|\leq\|A\|\|v\|.
\end{equation}
Thus $\|A^\times\|\leq\|A\|$. Now, by the previous lemma, for every $\overline u\in X$ there exists $\overline v\in Y^*$ such that $\|\overline v\|=1$ and $\overline v(A(\overline u))=\|A(\overline u)\|$. We now use the definition of the adjoint to get\[\begin{aligned}
\|A(\overline u)\|&=\overline v(A(\overline u))\\
&=(A^\times(\overline v))(\overline u)\\
&\leq\|A^\times(\overline v)\|\|\overline u\|\\
&\leq\|A^\times\|\|\overline v\|\|\overline u\|\\
&=\|A^\times\|\|\overline u\|.
\end{aligned}\]In the last equality we have used the fact that $\|\overline v\|=1$. Since $\overline v$ was arbitrary we now have $\|A\|\leq\|A^\times\|$, which completes the proof.
\end{Proof}
\endgroup
